\chapter{Conclusões e trabalhos futuros}
\label{cap:conclusao}

Esta dissertação investigou os gargalos de E/S da fase de simulação
final do algoritmo SDDP e propôs uma arquitetura paralela baseada em
escritas independentes com MPI-IO sobre o sistema de arquivos
Lustre, substituindo a estratégia atual com processo escritor único.
A avaliação experimental conduzida em dois ambientes distintos ---
AWS (FSx for Lustre dedicado) e Santos Dumont (LNCC, Lustre
compartilhado) --- confirma que a arquitetura proposta endereça
efetivamente o gargalo identificado e produz ganhos expressivos de
desempenho e escalabilidade.

Na configuração de 32 nós, a arquitetura MPI-IO reduz o tempo total
de execução em 36{,}7\% no Experimento AWS e em 40{,}4\% no Experimento
SDumont, em relação à implementação atual. A banda agregada
percebida pela aplicação cresce em 113$\times$ na AWS e
13,7$\times$ no SDumont. O contraste qualitativo entre as duas
implementações é tão relevante quanto a magnitude desses ganhos:
enquanto a implementação atual exibe colapso da banda a partir de
8~nós na AWS e de 16~nós no SDumont --- decorrente da serialização
inerente à implementação centralizada em um único processo escritor ---, a implementação proposta apresenta perfil
monotonicamente crescente de banda em toda a faixa avaliada,
indicando que a infraestrutura de armazenamento ainda não atingiu
seu limite nominal. Esse comportamento se sustenta em ambos os
ambientes, atestando a portabilidade da solução entre
infraestruturas heterogêneas --- nuvem dedicada e supercomputador
compartilhado.

A análise do impacto do número de \textit{Object Storage Targets}
(OSTs) reforça que o dimensionamento da camada de armazenamento é
parte essencial do ganho: passar de 4 para 10~OSTs no FSx for
Lustre amplia a banda agregada percebida em aproximadamente
5,6$\times$ em 16~nós e 23,6$\times$ em 32~nós, com redução do
tempo médio das escritas nas classes de bloco maiores. Esse
resultado sugere que, em \textit{deployments} futuros, o
\texttt{stripe\_count} e o \texttt{stripe\_size} do Lustre devem
constituir parâmetros de projeto, e não um detalhe
operacional.

\paragraph{Padrão de acesso dos arquivos diários.} Dos 117 arquivos
de resultado produzidos pelo SDDP nos experimentos, 104 seguem o
formato horário e 13 seguem o formato diário. Embora minoritários
em quantidade, os arquivos diários respondem por todas as escritas
de blocos muito pequenos observadas no \textit{workload}: cada
registro corresponde a um dia do mês e, diferentemente dos arquivos
horários --- nos quais a aplicação agrega os registros horários de um par
$(\text{estágio}, \text{cenário})$ em uma única escrita ---, ela não agrupa
os dias nos arquivos diários. A
sobreposição entre E/S e computação proporcionada pelas escritas
assíncronas mitiga parcialmente o impacto desse padrão fragmentado, mas a
frequência de pequenas requisições ao Lustre permanece como fonte residual
de sobrecarga e direciona uma das linhas de continuidade deste trabalho.

\paragraph{Disponibilidade de instâncias na nuvem.} A execução do
Experimento AWS enfrentou dificuldades de disponibilidade de
instâncias \texttt{r7i.12xlarge} na mesma região, dado o caráter
\textit{memory-optimized} dessa família e a demanda concorrente
sobre o \textit{pool} de recursos. Essa limitação operacional
restringiu a janela de experimentação, e o planejamento de futuras
execuções em nuvem para cargas HPC deve considerá-la.

\section*{Trabalhos futuros}

A partir dos resultados e limitações apresentados, este trabalho
identifica as seguintes direções de continuidade:

\begin{itemize}
  \item \textbf{Reorganização do padrão de acesso dos arquivos
        diários}, agrupando os registros mensais em uma única
        escrita por $(\text{estágio}, \text{cenário})$ ---
        analogamente ao que a aplicação já faz nos arquivos horários ---,
        de modo a eliminar a parcela residual de pequenas
        requisições ao Lustre identificada nestes experimentos.

  \item \textbf{Avaliação com adaptadores de rede de alta
        velocidade}, como InfiniBand HDR/NDR e \textit{Elastic
        Fabric Adapter} (EFA) da AWS, para investigar em que
        medida a rede deixa de ser fator limitante e como o custo
        da Coordenação MPI-IO escala em redes de baixa latência.

  \item \textbf{Sensibilidade ao tamanho de \textit{stripe} do
        Lustre}, com experimentos variando o \texttt{stripe\_size}
        para mapear o \textit{trade-off} entre fragmentação das
        requisições e distribuição efetiva entre os OSTs.

  \item \textbf{Buferização alinhada às \textit{stripes} do Lustre},
        substituindo o esquema atual --- em que cada \textit{rank}
        acumula em memória o bloco correspondente ao estágio
        completo de um cenário antes de submetê-lo em uma única
        chamada \texttt{MPI\_File\_iwrite\_at} ou
        \texttt{MPI\_File\_write\_at} --- por uma política que
        segmente as escritas pelos limites das \textit{stripes}, de
        modo que cada submissão MPI-IO incida sobre exatamente uma
        \textit{stripe} e, portanto, sobre exatamente um OST.
        Liao~\textit{et al.}~\cite{LiaoIPDPS:07} demonstram, em um
        contexto mais geral de escritas independentes em MPI-IO,
        que um esquema de bufferização em dois estágios alinhado ao
        particionamento do sistema de arquivos paralelo melhora
        significativamente o desempenho das escritas pequenas e
        desalinhadas; o mesmo princípio aplica-se diretamente ao
        perfil de E/S do SDDP horário, com potencial de reduzir a
        contenção entre \textit{ranks} concorrentes sobre os mesmos
        OSTs.

  \item \textbf{Sensibilidade ao limite síncrono/assíncrono}, com
        experimentos variando o ponto de corte entre escritas
        assíncronas e síncronas (por exemplo, 64, 128 e
        256~KB) para refinar empiricamente o parâmetro de projeto
        adotado neste trabalho.

  \item \textbf{Avaliação de \textit{two-phase I/O}}, isto é, de
        escritas coletivas com agregação interna pela biblioteca
        MPI-IO. Tentativas preliminares conduzidas durante esta
        pesquisa resultaram em \textit{deadlock}, atribuído à
        combinação de múltiplas escritas por cenário com o
        desalinhamento do número de cenários processados por cada
        \textit{rank} no padrão \textit{bag-of-tasks}. A
        viabilização dessa estratégia exigiria mecanismos de
        sincronização que reconciliem a natureza dinâmica da
        distribuição de tarefas com a semântica coletiva das
        operações de E/S.

  \item \textbf{Diagnóstico das variações no tempo de computação para
        uma mesma configuração}, caracterizando as flutuações
        observadas na fase de computação entre execuções com o mesmo
        número de nós. Como discutido na Seção~\ref{sec:sintese_avaliacoes},
        a eficiência paralela permanece abaixo do ideal linear no
        regime de \textit{strong scaling}: com o problema de tamanho
        fixo, custos aproximadamente constantes da computação ---
        sobretudo os associados ao acesso à memória --- deixam de ser
        amortizados à medida que se adicionam nós. A investigação
        detalhada desse gargalo computacional, distinto do gargalo de
        E/S tratado nesta dissertação (por exemplo, via
        \textit{profiling} da hierarquia de memória e da largura de
        banda de acesso), permitiria isolar seus efeitos dos de E/S e
        orientar otimizações futuras da fase de simulação.

  \item \textbf{Investigação da piora pontual observada em 4~nós no
        Experimento AWS} (Seção~\ref{sec:avaliacao_comparativa_aws}),
        em que o tempo total da implementação MPI-IO superou o da
        implementação atual apesar do ganho já mensurável de banda.
        Esse resultado é indício de que fatores externos ao módulo de
        E/S --- possivelmente variações no tempo de computação entre
        repetições --- dominaram esse ponto específico da curva. Sua
        caracterização, articulada com o diagnóstico das variações de
        computação mencionado acima, ajudaria a separar os efeitos de
        E/S dos efeitos computacionais na análise de escalabilidade.

  \item \textbf{Extensão da arquitetura à fase de geração de cortes do
        SDDP}, avaliando em que medida a estratégia de escritas
        independentes com MPI-IO se aplica ao padrão de comunicação e
        de sincronização dessa fase, distinto do da fase de simulação
        final tratada nesta dissertação.
\end{itemize}

A consolidação dessas direções pode contribuir para a construção de
uma arquitetura de E/S ainda mais escalável e portátil, aplicável
não apenas ao SDDP, mas a uma classe mais ampla de aplicações
científicas com padrões \textit{bag-of-tasks} e produção intensiva
de resultados.
